\documentclass[11pt]{amsart}
%\documentclass{article}
\usepackage[utf8]{inputenc}
\usepackage[T1]{fontenc}
\usepackage{amsmath, amssymb, amsthm}   
\usepackage{mathtools}
\usepackage{graphicx}
\usepackage[backend=biber,style=numeric]{biblatex} 
\usepackage{hyperref}
\usepackage{cleveref}

\addbibresource{refs.bib}

\title{Eigenmode Approximations of Noisy Tipping-Point Systems}
\author{Felix Shaw-Bell}
\address{Department of Mathematics, University of Auckland}
\email{fshawbell@gmail.com}
\subjclass[2020]{60H10, 60J60, 82C31, 37K55, 37H20}   % AMS Mathematics Subject Classification
\keywords{Fokker-Planck equation, stochastic dynamical systems, tipping points, perturbation methods}

\begin{document}
  \begin{abstract}
    Abstract goes here.
  \end{abstract}
  \maketitle
  \section{Introduction}
  \section{Setup and Preliminaries}
  \section{Results}
  \section{Discussion}
  \section{Conclusion}
  \section*{Code Availability}
    All code used to produce the results in this paper, including the eigenmode approximations and figure-generation scripts, is available at \url{https://github.com/fshawbell/eigenmode-tipping-point-systems}.
  \printbibliography
\end{document}